% GENERATED FILE. Edit canonical lessons, then run npm run generate:books.
% canonical-source-hash: fnv1a64:3032b47b684efbd2
% canonical-lessons: GU-C444-thad, GU-C444-chal, GU-C444-kali, GU-C444-jhumkhu, GU-C444-mahekvu

\chapter{Trunks, Bark and Buds}
\label{ch:gu-trunks-bark-and-buds}
\hlchaptermodality{444}

\begin{chapteropening}
\textbf{By the end of this chapter:} \emph{I can describe a tree from its trunk to its flowers.}

Complete the last lesson of chapter 444: I can describe a tree from its trunk to its flowers.
\end{chapteropening}

\section[\texorpdfstring{\gu{થડ} (thaḍ)}{થડ (thaḍ)}]{\gu{થડ} (thaḍ) — a tree trunk}

\label{lesson:GU-C444-thad}



\begin{quote}
Before the new one: say the Gujarati for a mango tree, then the Gujarati for to sprout.
\end{quote}

\subsection*{You'll want to know: \gu{થડ}}
\textbf{\gu{થડ}} — \emph{thaḍ} — \textquotedblleft{}a tree trunk\textquotedblright{}.

\begin{quote}
\gu{વડનું} \gu{થડ} \gu{ઘણું} \gu{મોટું} \gu{છે}.
\end{quote}

\begin{quote}
\emph{vaḍnũ thaḍ ghaṇũ moṭũ che.}
\end{quote}

\begin{quote}
The banyan's trunk is very big.
\end{quote}

A neuter noun.

\begin{cousinweb}
\textbf{\gu{થડ}} (\emph{thaḍ}) goes back to \emph{*thuḍa}, \textquotedblleft{}a trunk, a branch\textquotedblright{}, a form CDIAL reconstructs; CDIAL lists the Gujarati word (CDIAL 6098).
\end{cousinweb}

\begin{grammarlens}[title={What you've built}]
One more word for getting around.
\end{grammarlens}

\subsection*{Your turn}
\begin{itemize}
\raggedright
  \item \emph{Say it:} \emph{thaḍ}
  \item \emph{Say it:} \emph{thaḍ}, once more
  \item \emph{Say it:} \emph{vaḍnũ thaḍ ghaṇũ moṭũ che.} — The banyan's trunk is very big.
  \item \emph{Recall:} say the Gujarati for a mango tree, then the Gujarati for to sprout, then say \emph{thaḍ} again
\end{itemize}

\begin{tcolorbox}[breakable,colback=teal!4,colframe=teal!35!black,title={Before you move on}]
What does \gu{થડ} mean? (\textquotedblleft{}A tree trunk\textquotedblright{}.) Say it once more.
\end{tcolorbox}
\section[\texorpdfstring{\gu{છાલ} (chāl)}{છાલ (chāl)}]{\gu{છાલ} (chāl) — bark; peel}

\label{lesson:GU-C444-chal}



\begin{quote}
Before the new one: say the Gujarati for to sprout, then the Gujarati for a tree trunk.
\end{quote}

\subsection*{You'll want to know: \gu{છાલ}}
\textbf{\gu{છાલ}} — \emph{chāl} — \textquotedblleft{}bark; peel\textquotedblright{}.

\begin{quote}
\gu{આ} \gu{ઝાડની} \gu{છાલ} \gu{ઘણી} \gu{કઠણ} \gu{છે}.
\end{quote}

\begin{quote}
\emph{ā jhāḍnī chāl ghaṇī kaṭhaṇ che.}
\end{quote}

\begin{quote}
This tree's bark is very hard.
\end{quote}

A feminine noun.

\begin{cousinweb}
\textbf{\gu{છાલ}} (\emph{chāl}) is a tadbhava of Sanskrit \emph{challi}, \textquotedblleft{}bark\textquotedblright{}; CDIAL lists the Gujarati word, \textquotedblleft{}bark, peel\textquotedblright{} (CDIAL 5005). \gu{છોલવું}, to peel, is its neighbour.
\end{cousinweb}

\begin{grammarlens}[title={What you've built}]
One more word for getting around.
\end{grammarlens}

\subsection*{Your turn}
\begin{itemize}
\raggedright
  \item \emph{Say it:} \emph{chāl}
  \item \emph{Say it:} \emph{chāl}, once more
  \item \emph{Say it:} \emph{ā jhāḍnī chāl ghaṇī kaṭhaṇ che.} — This tree's bark is very hard.
  \item \emph{Recall:} say the Gujarati for to sprout, then the Gujarati for a tree trunk, then say \emph{chāl} again
\end{itemize}

\begin{tcolorbox}[breakable,colback=teal!4,colframe=teal!35!black,title={Before you move on}]
What does \gu{છાલ} mean? (\textquotedblleft{}Bark; peel\textquotedblright{}.) Say it once more.
\end{tcolorbox}
\section[\texorpdfstring{\gu{કળી} (kaḷī)}{કળી (kaḷī)}]{\gu{કળી} (kaḷī) — a bud}

\label{lesson:GU-C444-kali}



\begin{quote}
Before the new one: say the Gujarati for a tree trunk, then the Gujarati for bark.
\end{quote}

\subsection*{You'll want to know: \gu{કળી}}
\textbf{\gu{કળી}} — \emph{kaḷī} — \textquotedblleft{}a bud\textquotedblright{}.

\begin{quote}
\gu{ગુલાબના} \gu{છોડ} \gu{પર} \gu{નવી} \gu{કળીઓ} \gu{આવી}.
\end{quote}

\begin{quote}
\emph{gulābnā choḍ par navī kaḷīo āvī.}
\end{quote}

\begin{quote}
New buds came on the rose bush.
\end{quote}

A feminine noun.

\begin{cousinweb}
\textbf{\gu{કળી}} (\emph{kaḷī}) is a tadbhava of Sanskrit \emph{kali}, \textquotedblleft{}a bud\textquotedblright{}; CDIAL lists the Gujarati word (CDIAL 2934).
\end{cousinweb}

\begin{grammarlens}[title={What you've built}]
One more word for getting around.
\end{grammarlens}

\subsection*{Your turn}
\begin{itemize}
\raggedright
  \item \emph{Say it:} \emph{kaḷī}
  \item \emph{Say it:} \emph{kaḷī}, once more
  \item \emph{Say it:} \emph{gulābnā choḍ par navī kaḷīo āvī.} — New buds came on the rose bush.
  \item \emph{Recall:} say the Gujarati for a tree trunk, then the Gujarati for bark, then say \emph{kaḷī} again
\end{itemize}

\begin{tcolorbox}[breakable,colback=teal!4,colframe=teal!35!black,title={Before you move on}]
What does \gu{કળી} mean? (\textquotedblleft{}A bud\textquotedblright{}.) Say it once more.
\end{tcolorbox}
\section[\texorpdfstring{\gu{ઝૂમખું} (jhūmkhũ)}{ઝૂમખું (jhūmkhũ)}]{\gu{ઝૂમખું} (jhūmkhũ) — a bunch; a cluster}

\label{lesson:GU-C444-jhumkhu}



\begin{quote}
Before the new one: say the Gujarati for bark, then the Gujarati for a bud.
\end{quote}

\subsection*{You'll want to know: \gu{ઝૂમખું}}
\textbf{\gu{ઝૂમખું}} — \emph{jhūmkhũ} — \textquotedblleft{}a bunch; a cluster\textquotedblright{}.

\begin{quote}
\gu{આંબા} \gu{પર} \gu{કેરીનું} \gu{ઝૂમખું} \gu{લટકે} \gu{છે}.
\end{quote}

\begin{quote}
\emph{ā̃bā par kerīnũ jhūmkhũ laṭke che.}
\end{quote}

\begin{quote}
A bunch of mangoes hangs on the mango tree.
\end{quote}

A neuter noun.

\begin{cousinweb}
\textbf{\gu{ઝૂમખું}} (\emph{jhūmkhũ}) goes back to \emph{*jhuppa}, \textquotedblleft{}a tuft, a bunch\textquotedblright{}, a form CDIAL reconstructs; CDIAL lists the Gujarati word, \textquotedblleft{}a cluster\textquotedblright{} (CDIAL 5404).
\end{cousinweb}

\begin{grammarlens}[title={What you've built}]
One more word for getting around.
\end{grammarlens}

\subsection*{Your turn}
\begin{itemize}
\raggedright
  \item \emph{Say it:} \emph{jhūmkhũ}
  \item \emph{Say it:} \emph{jhūmkhũ}, once more
  \item \emph{Say it:} \emph{ā̃bā par kerīnũ jhūmkhũ laṭke che.} — A bunch of mangoes hangs on the mango tree.
  \item \emph{Recall:} say the Gujarati for bark, then the Gujarati for a bud, then say \emph{jhūmkhũ} again
\end{itemize}

\begin{tcolorbox}[breakable,colback=teal!4,colframe=teal!35!black,title={Before you move on}]
What does \gu{ઝૂમખું} mean? (\textquotedblleft{}A bunch; a cluster\textquotedblright{}.) Say it once more.
\end{tcolorbox}
\section[\texorpdfstring{\gu{મહેકવું} (mahekvũ)}{મહેકવું (mahekvũ)}]{\gu{મહેકવું} (mahekvũ) — to be fragrant; to smell sweet}

\label{lesson:GU-C444-mahekvu}



\begin{quote}
Before the new one: say the Gujarati for a bud, then the Gujarati for a cluster.
\end{quote}

\subsection*{You'll want to know: \gu{મહેકવું}}
\textbf{\gu{મહેકવું}} — \emph{mahekvũ} — \textquotedblleft{}to be fragrant; to smell sweet\textquotedblright{}.

\begin{quote}
\gu{સાંજે} \gu{બગીચામાં} \gu{ગુલાબનાં} \gu{ફૂલ} \gu{મહેકે} \gu{છે}.
\end{quote}

\begin{quote}
\emph{sā̃je bagīcāmā̃ gulābnā̃ phūl maheke che.}
\end{quote}

\begin{quote}
In the evening the roses smell sweet in the garden.
\end{quote}

Past (it …): \gu{તે} \gu{ફૂલ} \gu{મહેક્યું}, \emph{te phūl mahekyũ}.

\begin{cousinweb}
\textbf{\gu{મહેકવું}} (\emph{mahekvũ}) goes back to Sanskrit \emph{magha}, a word CDIAL takes, with some doubt, to have meant \textquotedblleft{}fragrance\textquotedblright{}; CDIAL lists the Gujarati verb, \textquotedblleft{}to give off a scent\textquotedblright{} (CDIAL 9703).
\end{cousinweb}

\begin{grammarlens}[title={What you've built}]
One more word for getting around.
\end{grammarlens}

\subsection*{Your turn}
\begin{itemize}
\raggedright
  \item \emph{Say it:} \emph{mahekvũ}
  \item \emph{Say it:} \emph{mahekvũ}, once more
  \item \emph{Say it:} \emph{sā̃je bagīcāmā̃ gulābnā̃ phūl maheke che.} — In the evening the roses smell sweet in the garden.
  \item \emph{Recall:} say the Gujarati for a bud, then the Gujarati for a cluster, then say \emph{mahekvũ} again
\end{itemize}

\begin{tcolorbox}[breakable,colback=teal!4,colframe=teal!35!black,title={Before you move on}]
What does \gu{મહેકવું} mean? (\textquotedblleft{}To be fragrant; to smell sweet\textquotedblright{}.) Say it once more.
\end{tcolorbox}
